\section{Strongly Rayleigh processes on countable groups}
\label{sec:strongly-rayleigh}

The finite response hypothesis of Proposition~\ref{prop:brownian-sampler}
holds for every strongly Rayleigh law. The fixed-cardinality covariance
identity for trees becomes an inequality for this larger class.
The covariance-row estimate below appears in Alishahi, Barzegar, and
Zamani~\cite[Lemma~2.5 and (1)]{AlishahiBarzegarZamani2023}, following the
nonnegative-row-sum result of Ghosh, Liggett, and
Pemantle~\cite[arXiv version, Lemma~3.2]{GhoshLiggettPemantle2017}.
We give a proof by symmetric homogenization and apply it uniformly after
positive external-field tilts.

\begin{lemma}[Strongly Rayleigh response]\label{lem:sr-response}
Let $F$ be a nonempty finite set and let $\mu$ be a strongly Rayleigh
probability law on $\{0,1\}^F$. For $h\in\R^F$, let
\[
 \mu_h(x)=\frac{\exp(\sum_{j\in F}h_jx_j)\mu(x)}
 {\sum_{z\in\{0,1\}^F}\exp(\sum_{j\in F}h_jz_j)\mu(z)},
 \qquad p_i(h)=\E_{\mu_h}X_i.
\]
For every $i,j\in F$, the derivative identity below holds, and for every
$i\in F$ the row satisfies
\begin{equation}\label{eq:sr-covariance-row}
 \frac{\partial p_i}{\partial h_j}
 =\Cov_{\mu_h}(X_i,X_j),
 \qquad
 \sum_{j\in F}\left|\frac{\partial p_i}{\partial h_j}\right|
 \le 2p_i(h)(1-p_i(h))\le\frac12.
\end{equation}
Consequently
\begin{equation}\label{eq:sr-response}
 |p_i(h)-p_i(h')|\le\frac12\max_{j\in F}|h_j-h'_j|
 \qquad(h,h'\in\R^F).
\end{equation}
\end{lemma}

\begin{proof}
The tilted generating polynomial is
\[
 g_{\mu_h}(z)
 =\frac{g_\mu((e^{h_j}z_j)_{j\in F})}
        {g_\mu((e^{h_j})_{j\in F})}.
\]
The denominator is positive. Positive coordinate scaling preserves
stability, so $\mu_h$ is strongly Rayleigh; see
Borcea, Br\"and\'en, and Liggett~\cite[Section~2.1(iv) and Proposition~3.1(2)]{BorceaBrandenLiggett2009}.
Differentiating the finite partition sum gives the covariance identity in
\eqref{eq:sr-covariance-row}.

Write $n=|F|$ and $\eta=\mu_h$. To recover a constant total count, take a
disjoint set $F'$ of $n$ dummy coordinates. First sample the occupied set
$A\subseteq F$ with law $\eta$; conditional on $A$, occupy a uniformly
chosen subset of $F'$ of size $n-|A|$. The resulting law
$\widetilde\eta$ keeps $\eta$ as its $F$-marginal and has exactly $n$
occupied coordinates in total. Write $\eta(A)$ for the probability of the
occupied set $A$, and let $e_k$ be the elementary symmetric polynomial
of degree $k$ in $w=(w_j)_{j\in F'}$. The generating polynomial of
$\widetilde\eta$, called the symmetric homogenization of $g_\eta$, is
\[
 \widetilde g(z,w)
 =\sum_{A\subseteq F}\eta(A)\prod_{j\in A}z_j\,
       \frac{e_{n-|A|}(w)}{\binom{n}{|A|}}.
\]
The polynomial has nonnegative coefficients, equals one at the
all-one vector, is homogeneous of degree $n$, and specializes to $g_\eta$
when all the $w$ variables equal one. The additional fact that it is stable
is supplied by
\cite[Definition~2.12 and Theorem~4.2]{BorceaBrandenLiggett2009}.
Thus $\widetilde\eta$ is strongly Rayleigh.

Let $C_{jk}=\Cov_{\widetilde\eta}(X_j,X_k)$. The Rayleigh inequality
in~\cite[Theorem~4.1]{BorceaBrandenLiggett2009}, evaluated at the all-one
vector, gives $C_{jk}\le0$ for $j\ne k$. Since the total count is constant,
for $i\in F$ we have
\[
 0=\sum_{j\in F\sqcup F'}C_{ij}
  =C_{ii}+\sum_{j\in F\setminus\{i\}}C_{ij}
            +\sum_{j\in F'}C_{ij}.
\]
The dummy covariances in the last sum are nonpositive. It follows that
\[
 -\sum_{j\in F\setminus\{i\}}C_{ij}
 =C_{ii}+\sum_{j\in F'}C_{ij}\le C_{ii},
 \qquad
 \sum_{j\in F}|C_{ij}|\le2C_{ii}.
\]
The covariances with both indices in $F$ are those of $\eta$, and
$C_{ii}=p_i(h)(1-p_i(h))\le1/4$. This proves
\eqref{eq:sr-covariance-row}, including degenerate coordinates. No
homogeneity or full-support assumption on $\mu$ was used. Integrating the
gradient along the segment from $h$ to $h'$ proves \eqref{eq:sr-response}.
\end{proof}

\begin{proof}[Proof of Theorem~\ref{thm:sr-group}]
Fix a site $i\in\Gamma$ and a finite observed set $S\subseteq\Gamma$.
Let $J=S\cup\{i\}$ and let $\nu_J$ be the marginal of $\nu$ on $J$.
For $t\ge0$ and $y\in\R^\Gamma$, tilt $\nu_J$ by
\[
 \exp\!\left(\sum_{j\in S}(y_j-t/2)x_j\right).
\]
The resulting probability law $\eta_{t,y}$ is an arbitrary positive
external-field tilt of the strongly Rayleigh law $\nu_J$, with zero log
field on $J\setminus S$. Its mean at $i$ is exactly the finite-tilt
function $b_{i,S}(t,y)$ in \eqref{eq:finite-tilt}. Finite differentiation
and Lemma~\ref{lem:sr-response} give
\[
 \frac{\partial b_{i,S}}{\partial y_j}(t,y)
 =\Cov_{\eta_{t,y}}(X_i,X_j)\quad(j\in S),
 \qquad
 \sum_{j\in S}\left|\frac{\partial b_{i,S}}{\partial y_j}(t,y)\right|
 \le\sum_{j\in J}|\Cov_{\eta_{t,y}}(X_i,X_j)|\le\frac12.
\]
Integrating along a segment of fields yields
\begin{equation}\label{eq:sr-finite-tilt}
 |b_{i,S}(t,y)-b_{i,S}(t,y')|
 \le\frac12\max_{j\in S}|y_j-y'_j|.
\end{equation}
For $S=\varnothing$, the finite-tilt function is constant and the same
statement holds with the empty maximum defined to be zero. The bound is
uniform in $i,S,t$ and in all finite field values. It uses the actual
finite marginal of $\nu$, with no infinite external-field tilt or
finite-volume approximation.

Choose finite sets $F_n\subseteq\Gamma$ increasing to $\Gamma$, with the
identity in every $F_n$, and put
\[
 S_n(h)=hF_n,\qquad h\in\Gamma.
\]
These are finite exhaustions at every site, and
$S_n(gh)=gS_n(h)$ for all $g,h\in\Gamma$. This requires neither finite
generation nor amenability; for a finite group one may take $F_n=\Gamma$.
Invariance of $\nu$ gives the exact finite-tilt identity
\begin{equation}\label{eq:sr-tilt-equivariance}
 b_{gh,gS}(t,\tau_g y)=b_{h,S}(t,y)
\end{equation}
for every $g,h,t,y$ and finite $S$. Indeed, in the numerator and denominator
of \eqref{eq:finite-tilt}, substitute $x=\tau_g z$. Then
$x_{gh}=z_h$, $(\tau_g y)_{gs}=y_s$, and $x_{gs}=z_s$, while the law of
$z$ is again $\nu$. All sums are finite, so the identity holds pointwise
even when the complete field is unbounded.

Define the limsup drift using these $S_n(h)$. Equations
\eqref{eq:sr-tilt-equivariance} and $S_n(gh)=gS_n(h)$ give
\[
 b_{gh}(t,\tau_g y)
 =\limsup_n b_{gh,S_n(gh)}(t,\tau_g y)
 =\limsup_n b_{h,S_n(h)}(t,y)=b_h(t,y).
\]
Thus the drift is equivariant under every regular left translation on its
entire domain. The uniform bound \eqref{eq:sr-finite-tilt} supplies the
hypothesis of Proposition~\ref{prop:brownian-sampler} with $L=1/2$.
That proposition gives a Borel map $\Psi$ from the product continuous-path
space to $\{0,1\}^{\Gamma}$, sending independent standard Brownian paths
to $\nu$ and commuting with every $\tau_g$ on every input.

The proof of Lemma~\ref{lem:edge-noise} constructs an everywhere-defined
Borel map from one uniform variable to a standard Brownian path, with fixed
defaults on its exceptional inputs. Applying this map coordinatewise
sends independent uniform labels on $\Gamma$ to independent Brownian paths
and commutes with all translations. Its composition with $\Psi$ is the
required map $\Phi_\nu$.
\end{proof}

The sets $F_n$ above need not be F\o lner sets. Their translated form is
available because the action is regular. For an arbitrary action on another
countable set, equivariant finite exhaustions are not automatic, and no
such extension is asserted here.

\begin{proof}[Proof of Corollary~\ref{cor:dpp-group}]
For a countable discrete set, every self-adjoint positive contraction
defines a determinantal probability law; see Lyons~\cite[Section~8]{Lyons2003}.
For a finite $J\subseteq\Gamma$, let $K_J$ be the compression of $K$ to
$\ell^2(J)$. It is a Hermitian positive contraction. For every $A\subseteq J$,
the inclusion probability of $A$ under the $J$-marginal is
\[
 \P(A\subseteq X)=\det K_{A,A}=\det (K_J)_{A,A}.
\]
These inclusion probabilities determine the finite binary law by
inclusion--exclusion. Hence that marginal is the determinantal law with
kernel $K_J$, and it is strongly Rayleigh by
\cite[Proposition~3.5]{BorceaBrandenLiggett2009}. This covers complex
Hermitian, singular, and projection kernels; no trace-class condition on
the operator on $\ell^2(\Gamma)$ is needed. Theorem~\ref{thm:sr-group}
now applies whenever the determinantal law is translation invariant.

If $K$ commutes with the left regular representation, its matrix entries
satisfy $K_{gh,gk}=K_{h,k}$. All finite inclusion determinants are then
translation invariant. Inclusion--exclusion and the uniqueness of a
countable law from its cylinder probabilities give invariance of the
determinantal law, proving the final assertion.
\end{proof}
